% !TEX root = ../main.tex

\ustcsetup{
  keywords  = {人形机器人, 全身运动控制, 遥操作, 强化学习, 感知控制闭环},
  keywords* = {Humanoid Robots, Whole-Body Control, Teleoperation,
               Reinforcement Learning, Perception-Control Loop},
}

\begin{abstract}
  人形机器人被认为是具身智能的重要载体，其在家庭服务、公共服务、灾害救援和高风险作业等场景中具有广阔应用前景。与传统机械臂或轮式移动平台相比，人形机器人需要在高维、强耦合和欠驱动条件下完成稳定移动与灵巧操作，因此全身运动控制一直是该领域的核心难题。围绕这一问题，本文以“人在回路”的研究范式为主线，尝试打通从人体动作采集、动作重定向、控制策略生成到感知驱动任务执行的完整技术链路，为人形机器人的通用全身运动控制提供一套可复用的实现框架。

  在数据采集层，本文构建了面向人形机器人遥操作的数据获取体系，综合利用 Noitom 惯性动捕、Nokov 光学动捕以及基于 SONIC 框架的仿真交互过程，覆盖从高自由度人体动作到复杂人机交互轨迹的多源异构数据。针对人体与机器人在骨骼长度、关节布局和自由度上的差异，进一步设计了基于 GMR 的运动重定向方法，将人体动作统一映射到 Unitree G1 及全尺寸机器人平台的可执行轨迹空间，为后续控制策略训练与真机部署提供稳定的数据基础。

  在控制策略层，本文提出了上下肢解耦的通用全身运动控制框架。下肢部分采用强化学习方法训练 Locomotion 策略，使机器人具备前进、后退、原地旋转和左右平移等基础运动能力，并通过域随机化和参数适配提升仿真到真实系统的迁移能力。上肢部分则以遥操作数据接口和逆运动学求解为核心，实现对上肢目标轨迹的高保真跟踪，并通过重心补偿和冲突处理机制完成全身协调控制，从而兼顾移动稳定性与操作精度。

  在任务验证层，本文开发了感知--控制闭环系统，并在两个代表性任务中验证所提方法的有效性。在 CyboRacket 动态球拍任务中，视觉系统对高速目标进行检测和轨迹预测，控制系统据此完成下肢导航、上肢挥拍和全身协同动作生成；在 OpenClaw 语音服务任务中，高层任务规划与视觉感知模块输出结构化子任务和目标位姿信息，控制框架进一步完成导航、抓取与递送等动作执行。实验与系统集成结果表明，本文所提出的技术路线能够有效支撑人形机器人从遥操作数据采集到自主任务执行的闭环应用。
\end{abstract}

\begin{abstract*}
  Humanoid robots are widely regarded as an important embodiment of embodied intelligence, with strong application potential in household service, public assistance, disaster response, and other complex scenarios. Compared with traditional robotic manipulators or wheeled platforms, humanoid systems must accomplish stable locomotion and dexterous manipulation under high-dimensional, strongly coupled, and underactuated conditions. Whole-body motion control therefore remains a central challenge. To address this problem, this thesis follows a human-in-the-loop paradigm and aims to establish an end-to-end technical pipeline spanning motion capture, motion retargeting, control policy generation, and perception-driven task execution.

  On the data acquisition side, a multi-source teleoperation data collection system is constructed by integrating Noitom inertial motion capture, Nokov optical motion capture, and simulation-based interaction data generated within the SONIC framework. This design covers both high-degree-of-freedom human motion and complex loco-manipulation trajectories. To bridge the morphological gap between humans and humanoid robots, including differences in limb length, joint layout, and actuation space, a GMR-based motion retargeting method is further introduced to map human motion into executable trajectories for the Unitree G1 platform and full-scale humanoid robots.

  On the control side, this thesis proposes a decoupled whole-body control framework in which the lower body focuses on locomotion stability and the upper body focuses on manipulation accuracy. The lower body is driven by reinforcement learning policies that support forward motion, backward motion, in-place turning, and lateral translation, while domain randomization and parameter adaptation are employed to improve sim-to-real transfer. The upper body relies on teleoperation interfaces and inverse kinematics solvers to track target trajectories with high fidelity. Whole-body coordination is achieved through conflict handling and center-of-mass compensation mechanisms, enabling stable and flexible combined motions.

  On the task validation side, a perception-control closed-loop system is developed and evaluated in two representative scenarios. In the CyboRacket dynamic racket task, the vision module detects and predicts high-speed targets, and the proposed controller generates coordinated lower-body navigation and upper-body striking motions. In the OpenClaw voice-commanded service task, high-level planning and visual perception provide structured subtasks and target poses, while the whole-body controller executes navigation, grasping, and delivery behaviors. The resulting system demonstrates that the proposed pipeline can effectively support humanoid robots from teleoperation data acquisition to autonomous task execution.
\end{abstract*}
